\begin{tikzpicture}[font=\small,>=stealth,line join=round]
\draw[gray!40,rounded corners=3pt] (0,0) rectangle (14.2,10.2);
\node at (7.1,9.6) {单位圆与平行直线相切：两个边界都避开缺点 $E$};
\begin{scope}[shift={(4.0,4.6)},x=2cm,y=2cm]
\draw[->,gray] (-1.3,0)--(1.45,0) node[right]{$u$};
\draw[->,gray] (0,-1.4)--(0,1.6) node[above]{$v$};
\draw[blue,very thick] (0,0) circle(1);
\draw[teal!80!black,thick] (-1.3,{sqrt(5)/2+.65})--(1.3,{sqrt(5)/2-.65});
\draw[orange!85!black,thick] (-1.3,{-sqrt(5)/2+.65})--(1.3,{-sqrt(5)/2-.65});
\fill[teal!80!black] ({1/sqrt(5)},{2/sqrt(5)}) circle(2pt) node[above right=5pt]{$A$};
\fill[orange!85!black] ({-1/sqrt(5)},{-2/sqrt(5)}) circle(2pt) node[below left=5pt]{$B$};
\draw[red!70!black,thick,fill=white] (1,0) circle(2.8pt) node[below right=6pt]{$E$};
\node[below left=4pt] at (0,0){$O$};
\end{scope}
\node[align=left,anchor=west] at (8.1,5.6){$u^2+v^2=1$\\[10pt]$\tfrac12u+v=\tfrac32-y$\\[12pt]\textcolor{teal!80!black}{$A$ 对应最小值 $L$}\\[9pt]\textcolor{orange!85!black}{$B$ 对应最大值 $U$}\\[12pt]\textcolor{red!70!black}{$E=(1,0)$ 不可达}\\[7pt]两切点都不是 $E$};
\node[text width=13cm,align=center] at (7.1,.6) {圆的横纵坐标等比例；空心点不属于原参数能够到达的轨迹。};
\end{tikzpicture}
